\section{部署与运行方案}

本项目部署与运行方案以“课程验收可复现、核心链路可演示、配置边界清晰、第三方异常可降级”为原则。本期交付采用 Vue 3 + Vite 前端、FastAPI 后端、根目录 \texttt{data/} JSON 集合和高德/DeepSeek 可选第三方服务，不使用 Docker，也不在本期强制引入 PostgreSQL、Redis、对象存储或正式网关。

\subsection{运行环境}

课程验收环境划分为本地开发环境和演示测试环境两类。

\begin{itemize}
  \item \textbf{本地开发环境：}前端通过 \texttt{npm install} 和 \texttt{npm run dev} 启动，Vite 默认端口为 5173；后端在 conda 环境 \texttt{fish} 中进入 \texttt{backend} 目录，通过 \texttt{uvicorn app.main:app --reload --host 127.0.0.1 --port 8000} 启动。
  \item \textbf{演示测试环境：}可在同一主机或测试服务器上运行前端构建产物与 FastAPI 服务，使用统一的环境变量文件配置高德和 DeepSeek Key。若未配置真实 Key，系统应使用内置演示数据或 mock 解释保持主流程可演示。
\end{itemize}

当前 MVP 阶段根目录 \texttt{data/} 使用 JSON 模拟 NoSQL 数据库，每个 \texttt{.json} 对应一个集合，每个 \texttt{.schema.json} 对应集合字段说明。JSON 文件适合课堂 Demo 和单机演示，不承诺正式多人并发写入能力。

\subsection{本期部署拓扑}

本期部署拓扑为：浏览器访问 Vue 3 + Vite 前端页面；前端通过 \texttt{api.js} 访问 \texttt{/api}；Vite 开发代理或测试环境反向代理将请求转发到 FastAPI；FastAPI 路由调用 Service 层；Service 层读取或写入 JSON 集合，并按需调用高德或 DeepSeek 服务。推荐排序由后端规则模型完成，AI 只负责解释生成，不直接改变排序。

第三方 Key 均通过后端环境变量配置，不允许出现在 Git 仓库、前端打包文件、公开截图和测试报告中。高德 JS API 配置由 \texttt{/api/amap/config} 下发；高德 Web Service Key 和 DeepSeek Key 只保存在后端环境变量中。

\subsection{配置项与启动检查}

后端至少需要检查以下配置项，配置项清单见表 \ref{tab:config-items}。

\begin{table}[htbp]
  \centering\small
  \caption{核心配置项清单}
  \label{tab:config-items}
  \begin{tabularx}{\textwidth}{>{\raggedright\arraybackslash}p{2.8cm}|>{\raggedright\arraybackslash}p{1.5cm}|X|>{\raggedright\arraybackslash}p{2.8cm}|>{\raggedright\arraybackslash}p{1.2cm}}
    \toprule
    \textbf{配置项} & \textbf{类型} & \textbf{说明} & \textbf{默认值/示例} & \textbf{必须} \\
    \midrule
    \texttt{APP\_ENV} & 字符串 & 运行环境，决定 AI 和第三方服务是否启用真实调用 & \texttt{development} & 是 \\
    \hline
    \texttt{BACKEND\_HOST} & 字符串 & 后端监听地址 & \texttt{127.0.0.1} & 是 \\
    \hline
    \texttt{BACKEND\_PORT} & 整数 & 后端监听端口 & \texttt{8000} & 是 \\
    \hline
    \texttt{AMAP\_JS\_API\_KEY} & 字符串 & 前端加载高德 JS API 所需 Key & 高德注册获取 & 是 \\
    \hline
    \texttt{AMAP\_SECURITY\_CODE} & 字符串 & 高德 JS API 安全密钥 & 高德注册获取 & 否 \\
    \hline
    \texttt{AMAP\_WEB\_SERVICE\_KEY} & 字符串 & 后端调用高德 Web Service 的 Key & 高德注册获取 & 否 \\
    \hline
    \texttt{DEEPSEEK\_API\_KEY} & 字符串 & AI 解释服务 API Key，可为空 & 空（启用 mock） & 否 \\
    \hline
    \texttt{DEEPSEEK\_MODEL} & 字符串 & AI 解释服务模型名称 & \texttt{deepseek-chat} & 否 \\
    \hline
    \texttt{DATA\_DIR} & 路径 & JSON 数据目录，指向根目录 \texttt{data/} & \texttt{../data} & 是 \\
    \hline
    \texttt{LOG\_LEVEL} & 字符串 & 日志级别 & \texttt{INFO} & 否 \\
    \bottomrule
  \end{tabularx}
\end{table}

正式提交部署材料时，应提供不含真实密钥的环境变量模板；真实 \texttt{.env} 由部署人员在本机或服务器上配置，不进入版本控制。

\subsection{降级与容错设计}

系统需要具备模块化降级能力，各模块降级策略见表 \ref{tab:degradation}。

\begin{table}[htbp]
  \centering\small
  \caption{模块降级策略表}
  \label{tab:degradation}
  \begin{tabularx}{\textwidth}{>{\raggedright\arraybackslash}p{2.2cm}|X|>{\raggedright\arraybackslash}p{3.0cm}|>{\raggedright\arraybackslash}p{1.6cm}}
    \toprule
    \textbf{故障场景} & \textbf{降级行为} & \textbf{用户感知} & \textbf{影响范围} \\
    \midrule
    高德地图底图加载失败 & 展示钓点列表和错误提示，隐藏地图容器 & 页面显示文字提示，仍可浏览钓点列表 & 地图页 \\
    \hline
    天气接口超时或失败 & 采用最近一次天气快照；若无缓存则展示”天气未知” & 天气模块显示旧数据或占位符 & 推荐、记录、详情页 \\
    \hline
    AI 解释接口失败 & 推荐排序结果正常返回，解释部分替换为规则模板或”暂未生成” & 推荐列表正常，解释文字较简略 & 推荐模块 \\
    \hline
    高德 POI 搜索不可用 & 继续使用 \texttt{data/pois.json} 中的本地演示钓点 & 搜索结果仅包含本地数据 & POI 搜索 \\
    \hline
    DeepSeek Key 未配置 & 启用 mock 模式，返回预设解释模板 & 无感知，解释内容为通用模板 & AI 解释 \\
    \hline
    JSON 数据文件写入失败 & 前端提示”保存失败，请重试”，保留表单数据不清空 & 弹窗提示，可手动重试 & 记录、发布 \\
    \bottomrule
  \end{tabularx}
\end{table}

活动、订单、群聊和会员报告属于本期可演示扩展模块。若现场时间或网络条件不足，应优先展示地图、推荐、记录、报告和社区五个主链路，避免为了展示扩展模块影响核心验收。

\subsection{正式工程化部署规划}

正式多人使用时，可在本期架构基础上逐步引入 PostgreSQL/PostGIS、Redis、对象存储、HTTPS 域名、Nginx 反向代理、统一鉴权、监控告警和备份恢复。该阶段的目标是支持多人并发写入、图片视频管理、权限控制和长期运行，不属于本期课程 MVP 的验收承诺。

正式化迁移应遵循向后兼容原则：先保持 API 返回结构稳定，再迁移数据层；新增字段优先设置默认值，不在同一版本中删除旧字段；前端接口调用先保留兼容层，再逐步切换到新字段。通过这种方式，课程 Demo 可以平滑演进为后续工程版本。

\subsection{交付与运行支撑}

本期运行支撑材料应包括部署说明、系统操作手册、接口说明、JSON 集合字段说明、测试账号或演示用户、演示数据和常见问题清单。课堂验收时，可采用本地单机部署或测试服务器演示两种方式：本地部署便于老师检查项目代码与架构，在线演示便于答辩展示真实交互。

运行维护期间，平台支持按模块降级：真实支付未接入时，装备订单仅保留模拟下单与取消；AI 接口费用超限时，推荐排序继续运行，仅暂停解释生成；活动服务未开放时，地图、记录、社区和报告功能仍可独立运行。通过模块化部署和降级策略，系统能够在课程验收、演示竞赛和后续迭代之间保持一致的技术路线。
